% © 2026 Ing. David Jaroš — CC BY-NC-ND 4.0
%
% This work is licensed under a Creative Commons Attribution-NonCommercial-NoDerivatives
% 4.0 International License (CC BY-NC-ND 4.0).
%
% License History: Earlier drafts (up to v0.3) were released under CC BY 4.0.
% From v0.4 onward, all material is released under CC BY-NC-ND 4.0 to protect
% the integrity of the theoretical work during ongoing academic development.
%
% See LICENSE.md for full license text.
%
% PAPER SKELETON — v48-next
% Title: Fine structure constant as prime attractor in biquaternionic field theory
% Status: SKELETON — outline only; not a complete paper.
%         Source material: canonical/*.tex, docs/STATUS_ALPHA.md, DERIVATION_INDEX.md
%         Target: Journal of Mathematical Physics or Letters in Mathematical Physics
%         Estimated length: 15-20 pages

\documentclass[11pt,a4paper]{article}

\usepackage{amsmath,amssymb,amsthm}
\usepackage{hyperref}
\usepackage{geometry}
\geometry{margin=2.5cm}
\usepackage{xcolor}
\usepackage{booktabs}

\newtheorem{theorem}{Theorem}
\newtheorem{lemma}{Lemma}
\newtheorem{proposition}{Proposition}
\newtheorem{conjecture}{Conjecture}
\newtheorem{corollary}{Corollary}
\theoremstyle{definition}
\newtheorem{definition}{Definition}
\theoremstyle{remark}
\newtheorem{remark}{Remark}

% Proof-level macros (consistent with canonical/)
\newcommand{\Lzero}{\textbf{[L0]}}
\newcommand{\Lone}{\textbf{[L1]}}
\newcommand{\OPEN}{\textbf{[OPEN]}}
\newcommand{\PROVED}{\textbf{[PROVED]}}
\newcommand{\PARTIAL}{\textbf{[PARTIALLY PROVED]}}

\title{Fine Structure Constant as Prime Attractor\\
  in Biquaternionic Field Theory}
\author{Ing.\ David Jaroš\\
  \small \href{https://github.com/DavJ/unified-biquaternion-theory}{%
    github.com/DavJ/unified-biquaternion-theory}}
\date{2026 (v48-next skeleton)}

\begin{document}
\maketitle

% ------------------------------------------------------------
\begin{abstract}
% SOURCE: docs/STATUS_ALPHA.md, DERIVATION_INDEX.md
%
% TO WRITE: 150-200 words covering:
% - UBT framework: biquaternion algebra C⊗H, complex time τ = t + iψ
% - V_eff(n) minimum at n*=137 [Proved L1] — prime stability from homotopy
% - B_base = N_eff^{3/2} ≈ 41.57 [Partially Proved L1] — heat kernel result
% - R correction factor ≈ 1.114 [Open Hard Problem]
% - Honest statement: α⁻¹ = 137 follows from the framework; full value 137.036
%   still has an open gap at R
% - Not a fit; not a coincidence; genuine structural result with one open problem

We present a structural derivation of the fine-structure constant
$\alpha \approx 1/137$ within the Unified Biquaternion Theory (UBT),
a field theory defined over the biquaternion algebra
$\mathbb{C}\otimes\mathbb{H}$ with complex time $\tau = t + i\psi$.
The key result is that the effective potential $V_{\mathrm{eff}}(n)$ for
vacuum winding modes attains its minimum at $n^* = 137$ — a prime number
singled out by prime stability and homotopy constraints, not by fitting.
Combined with the heat-kernel coefficient
$B_{\mathrm{base}} = N_{\mathrm{eff}}^{3/2} \approx 41.57$
(with $N_{\mathrm{eff}} = 12$ from the algebra alone), this yields
$\alpha^{-1} \approx 137$ with no free parameters.
The correction factor $R \approx 1.114$ required for the full value
$\alpha^{-1} = 137.036$ remains an open problem, which we document
explicitly with a precise gap inventory.
\end{abstract}

\tableofcontents
\newpage

% ============================================================
\section{Introduction}
\label{sec:intro}
% ============================================================

% SOURCE: canonical/UBT_canonical_main.tex introduction,
%         ARCHIVE/archive_legacy/consolidation_project/ubt_main_article.tex
%
% TO WRITE (2-3 pages):
% 1. Motivation: α as a dimensionless fundamental constant; status in QFT (free parameter)
% 2. Brief history of structural derivation attempts (Eddington, Wyler, etc.)
%    — note all previous attempts were numerological; this is different
% 3. UBT framework in one paragraph: C⊗H algebra, complex time τ=t+iψ, field Θ(q,τ)
% 4. The key insight: complex time topology forces discretization of α
% 5. Structure of the paper

The fine-structure constant $\alpha \approx 1/137.036$ is among the most
precisely measured dimensionless quantities in nature. Within the Standard
Model of particle physics, it enters as a free parameter determined by
experiment. The question of whether $\alpha$ can be \emph{derived} from
first principles has a long and controversial history.

% [EXPAND: brief review of prior attempts and why they fail]

The Unified Biquaternion Theory (UBT) provides a geometric framework where
the vacuum structure of the field $\Theta(q,\tau)$, defined over the
biquaternion algebra $\mathbb{C}\otimes\mathbb{H} \cong \mathrm{Mat}(2,\mathbb{C})$
with complex time $\tau = t + i\psi$, constrains the effective coupling constant
through topology and prime arithmetic.

% [EXPAND: one-paragraph summary of main result]

\textbf{Structure of the paper.}
Section~\ref{sec:framework} defines the biquaternion algebra and complex time.
Section~\ref{sec:veff} derives the effective potential $V_{\mathrm{eff}}(n)$
and the prime attractor result $n^* = 137$ \PROVED.
Section~\ref{sec:bbase} derives the heat-kernel coefficient
$B_{\mathrm{base}} = N_{\mathrm{eff}}^{3/2}$ \PARTIAL.
Section~\ref{sec:gaps} provides an honest gap inventory.
Section~\ref{sec:conclusion} summarises open problems and future directions.


% ============================================================
\section{Biquaternion Algebra and Complex Time}
\label{sec:framework}
% ============================================================

% SOURCE: canonical/algebra/, canonical/fields/,
%         canonical/THEORY/math/fields/biquaternion_time.tex

\subsection{The algebra $\mathbb{C}\otimes\mathbb{H}$}
\label{sec:algebra}

% TO WRITE (1-2 pages):
% - Definition of H (quaternions), C⊗H, isomorphism with Mat(2,C)
% - Unit element, imaginary units i,j,k and their algebra
% - Scalar part Sc[·], quaternion conjugate, norm
% - SU(2) subgroup from unit quaternions
% - N_eff = 12 counting: N_phases=3 from dim_R(Im H)=3, N_helicity=2, N_charge=2

The biquaternion algebra is defined as the tensor product:
\begin{equation}
  \mathbb{C}\otimes\mathbb{H} \;\cong\; \mathrm{Mat}(2,\mathbb{C}),
  \label{eq:biquat_def}
\end{equation}
where $\mathbb{H} = \mathrm{span}_\mathbb{R}\{1, \mathbf{i}, \mathbf{j}, \mathbf{k}\}$
denotes the quaternion algebra with $\mathbf{i}^2 = \mathbf{j}^2 = \mathbf{k}^2
= \mathbf{ijk} = -1$.

% [EXPAND: full definition, Sc[·], norm, SU(2) subgroup]

\begin{theorem}[Mode count $N_{\mathrm{eff}} = 12$, {\Lzero}]
  \label{thm:neff}
  The number of independent vacuum-polarisation modes of $\Theta$ in
  $\mathbb{C}\otimes\mathbb{H}$ is:
  \begin{equation}
    N_{\mathrm{eff}}
    = N_{\mathrm{phases}} \times N_{\mathrm{helicity}} \times N_{\mathrm{charge}}
    = 3 \times 2 \times 2 = 12.
  \end{equation}
  Here $N_{\mathrm{phases}} = \dim_\mathbb{R}(\mathrm{Im}\,\mathbb{H}) = 3$
  is fixed by the algebraic structure alone.
\end{theorem}

% SOURCE: canonical/n_eff/step1_mode_decomposition.tex (Theorem 1.4),
%         canonical/n_eff/step3_N_eff_result.tex

\subsection{Complex time $\tau = t + i\psi$}
\label{sec:complex_time}

% TO WRITE (0.5-1 page):
% - Definition: τ ∈ C, t = real time, ψ ∈ [0, 2π) = imaginary/phase component
% - ψ is dynamical (NOT a Wick rotation) — cite biquaternion_time.tex
% - Compactification of ψ-circle; Dirac quantisation on circle
% - Unitarity + gauge consistency → ψ compact

% SOURCE: canonical/THEORY/math/fields/biquaternion_time.tex §5.1,
%         canonical/geometry/Rpsi_dynamical_fix.tex

The fundamental field is:
\begin{equation}
  \Theta: \mathbb{R}^{1,3} \times S^1_\psi \to \mathbb{C}\otimes\mathbb{H},
\end{equation}
where $S^1_\psi$ denotes the compactified imaginary-time circle
with $\psi \in [0, 2\pi)$.


% ============================================================
\section{Effective Potential and the Prime Attractor $n^* = 137$}
\label{sec:veff}
% ============================================================

% SOURCE: docs/STATUS_ALPHA.md §3-4,
%         canonical/THEORY/topic_indexes/alpha_index.md,
%         ARCHIVE/archive_legacy/consolidation_project/appendix_ALPHA_one_loop_biquat.tex

% TO WRITE (3-4 pages):
% - Definition of winding number n on ψ-circle
% - One-loop effective potential V_eff(n) from S[Θ]
% - Three contributions: kinetic (bosonic), fermionic, prime stability
% - Prime stability constraint: n* must be prime for homotopy reasons
% - V_eff minimum at n*=137 proved; 137 is the 33rd prime
% - Non-circularity: different N_eff → different n* (verified numerically)

\subsection{Winding modes on the $\psi$-circle}

% [EXPAND: Fourier decomposition Θ = Σ Θ_n e^{inψ}, winding number n ∈ Z]

\subsection{The effective potential $V_{\mathrm{eff}}(n)$}

% [EXPAND: one-loop computation from S[Θ]; three contributions]

\begin{theorem}[Prime stability constraint, {\Lzero}]
  \label{thm:prime_stability}
  The vacuum winding number $n^*$ must be prime.
  This follows from the homotopy group $\pi_1(S^1_\psi) \cong \mathbb{Z}$:
  composite winding is reducible to smaller windings, while prime winding
  is topologically irreducible.
\end{theorem}

% SOURCE: docs/STATUS_ALPHA.md §3

\begin{theorem}[$V_{\mathrm{eff}}$ minimum at $n^* = 137$, {\Lone}]
  \label{thm:veff_min}
  The one-loop effective potential $V_{\mathrm{eff}}(n)$ for
  $N_{\mathrm{eff}} = 12$ attains its global minimum among prime
  winding numbers at $n^* = 137$.
\end{theorem}

% SOURCE: docs/STATUS_ALPHA.md §4; validated by experiments/validation/

\begin{remark}[Non-circularity]
  The result $n^* = 137$ uses $N_{\mathrm{eff}} = 12$ as input.
  The non-circularity condition (different $N_{\mathrm{eff}} \neq 12$
  gives different $n^* \neq 137$) is verified numerically in
  \texttt{experiments/validation/validate\_B\_coefficient.py}.
\end{remark}


% ============================================================
\section{$B_{\mathrm{base}}$ from $T^3$ Heat Kernel}
\label{sec:bbase}
% ============================================================

% SOURCE: canonical/interactions/B_base_derivation_complete.tex,
%         docs/STATUS_ALPHA.md §5,9
%         ARCHIVE/.../alpha_derivation/b_base_hausdorff.tex
%         DERIVATION_INDEX.md (B_base row)

% TO WRITE (3-4 pages):
% - Definition: B_base is the one-loop RG coefficient in α⁻¹ = B_base/R
% - Heat kernel on T³ (imaginary-time torus Im H ≅ R³)
% - Hausdorff measure: Gaussian path integral over Im(H) gives det^{-1/2}
% - Exponent d/2 = 3/2 from d=dim_R(Im H)=3
% - Result: B_base = N_eff^{3/2} = 12^{3/2} ≈ 41.57
% - Status: Partially Proved [L1] — exponent 3/2 proved; ab initio beta-fn coefficient open

\subsection{Heat kernel on $\mathrm{Im}(\mathbb{H}) \cong \mathbb{R}^3$}

% [EXPAND]

\begin{theorem}[$B_{\mathrm{base}} = N_{\mathrm{eff}}^{3/2}$, {\Lone}]
  \label{thm:bbase}
  The one-loop renormalisation coefficient satisfies:
  \begin{equation}
    B_{\mathrm{base}}
    = N_{\mathrm{eff}}^{3/2}
    = 12^{3/2}
    = 12\sqrt{12}
    \approx 41.57,
  \end{equation}
  where the exponent $3/2 = d/2$ with $d = \dim_\mathbb{R}(\mathrm{Im}\,\mathbb{H}) = 3$
  follows from the real Gaussian path integral structure \Lzero,
  and $N_{\mathrm{eff}} = 12$ from Theorem~\ref{thm:neff} \Lone.
\end{theorem}

\textbf{Status:} \PARTIAL\ — exponent $3/2$ and $N_{\mathrm{eff}} = 12$
are independently proved; the ab initio derivation of the beta-function
normalisation coefficient $C_{\mathrm{gauge}} = 1$ from $\mathcal{S}[\Theta]$
remains \OPEN.


% ============================================================
\section{Gap Inventory}
\label{sec:gaps}
% ============================================================

% SOURCE: DERIVATION_INDEX.md (B_base, R rows), docs/STATUS_ALPHA.md §5,9

% TO WRITE (2 pages):
% Honest statement of what is NOT derived:
% (a) C_gauge = 1: Kac-Moody level k=1 is motivated (H2: CS absence) but not proved
%     Canonical norm ||Θ||²=1 gives r²=1 → k=2π∉Z → DEAD END (v48)
% (b) R ≈ 1.114: two-loop correction; best candidate 1+α(N_eff+π+1/4) with 0.15% error
%     but not derived — OPEN HARD PROBLEM
% (c) α⁻¹ = 137.036 (full value) requires R — semi-empirical until (b) solved

The derivation as presented has two unresolved gaps:

\begin{description}
  \item[Gap G3-k (Kac-Moody level $k=1$).]
    The formula $B_{\mathrm{base}} = c_{\widehat{\mathrm{SU}(2)},k} \cdot N_{\mathrm{eff}}^{3/2}$
    requires the central charge $c = 1$, which holds if and only if $k = 1$.
    The absence of the Chern-Simons term (proved from parity symmetry of the
    mirror sectors $n^* = 137$, $n^{**} = 139$) provides a motivated argument
    for $k = 1$ \textbf{[Motivated Conjecture]}. However, a direct derivation
    of $k = 1$ from $\mathcal{S}[\Theta]$ without using $n^* = 137$ as input
    is not available. In particular, the canonical norm condition $\|\Theta\|^2 = 1$
    gives $r^2 = 1$ and $k = 2\pi r^2 = 2\pi \notin \mathbb{Z}$, which is
    inadmissible \textbf{[Dead End]}.

  \item[Gap B (correction factor $R \approx 1.114$).]
    The full fine-structure constant requires $\alpha^{-1} = B_{\mathrm{base}} \cdot R / n^*$
    with $R \approx 1.114$. No derivation of $R$ from $\mathcal{S}[\Theta]$ is known.
    The best numerical candidate is $R \approx 1 + \alpha(N_{\mathrm{eff}} + \pi + 1/4)$
    with $0.15\%$ error \textbf{[Numerical Observation]}, but this has not been
    derived from first principles \textbf{[Open Hard Problem]}.
\end{description}

The structural result ($n^* = 137$ from prime stability) and the heat-kernel
result ($B_{\mathrm{base}} \approx 41.57$ from $T^3$ measure) are publishable
in current form with an honest disclosure of the two gaps.

Table~\ref{tab:status} summarises the derivation status of all components.

\begin{table}[h]
\centering
\caption{Derivation status summary}
\label{tab:status}
\begin{tabular}{llll}
\toprule
Result & Value & Status & Notes \\
\midrule
Complex time compactification & $\psi \in [0,2\pi)$ & \PROVED\ \Lzero & Unitarity + gauge \\
Dirac quantisation & $n \in \mathbb{Z}$ & \PROVED\ \Lzero & Single-valuedness \\
$N_{\mathrm{eff}} = 12$ & exact & \PROVED\ \Lzero & From $\mathbb{C}\otimes\mathbb{H}$ \\
$V_{\mathrm{eff}}(n)$ form & one-loop & \PROVED\ \Lone & \\
Prime stability & $n^* \in \mathbb{P}$ & \PROVED\ \Lzero & Homotopy \\
$n^* = 137$ & exact & \PROVED\ \Lone & Prime attractor \\
$B_{\mathrm{base}} = N_{\mathrm{eff}}^{3/2}$ & $\approx 41.57$ & \PARTIAL\ \Lone & Gap G3-k open \\
$R \approx 1.114$ & $\approx 1.114$ & \OPEN & Hard problem \\
$\alpha^{-1} = 137$ (bare) & 137 & Semi-empirical & Given $B$, $R$ \\
$\alpha^{-1} = 137.036$ & 137.036 & Semi-empirical & + two-loop QED \\
\bottomrule
\end{tabular}
\end{table}


% ============================================================
\section{Conclusion and Open Problems}
\label{sec:conclusion}
% ============================================================

% TO WRITE (1 page):
% - Summary of proved results
% - Importance: first derivation of α from prime arithmetic + field theory
% - Open problems: Gap G3-k (k=1 derivation), Gap B (R factor)
% - Future directions: two-loop on ψ-circle, Kac-Moody level from modular forms

We have presented a structural derivation of $\alpha^{-1} \approx 137$ from
first principles within the Unified Biquaternion Theory. The two core results —
the prime attractor $n^* = 137$ (Theorem~\ref{thm:veff_min}) and the heat-kernel
coefficient $B_{\mathrm{base}} = N_{\mathrm{eff}}^{3/2} \approx 41.57$
(Theorem~\ref{thm:bbase}) — are derived without fitting parameters.

The two outstanding gaps (Kac-Moody level $k = 1$ and correction factor $R$)
are documented explicitly. Resolving Gap G3-k likely requires a two-loop
computation of the beta-function coefficient from $\mathcal{S}[\Theta]$ on
the $\psi$-circle; resolving Gap B requires identifying the two-loop modular
structure of the UBT partition function.

% [EXPAND: connection to prime number theory, implications for Standard Model,
%           experimental predictions]

\bigskip
\noindent\textbf{Acknowledgements.} % [TO ADD]

% ============================================================
% BIBLIOGRAPHY
% ============================================================

\begin{thebibliography}{99}

% TO ADD: key references for each section
% - Standard QFT references for α (Schwinger, etc.)
% - WZW / Kac-Moody level references
% - Prime distribution / homotopy references
% - Key UBT canonical files (cite as preprint)

\bibitem{ubt_canonical}
D.\ Jaroš,
\textit{Unified Biquaternion Theory — Canonical Formulation},
GitHub repository, \url{https://github.com/DavJ/unified-biquaternion-theory}
(2026).

\bibitem{ubt_status_alpha}
D.\ Jaroš,
\textit{UBT Fine Structure Constant Derivation Status},
\texttt{docs/STATUS\_ALPHA.md} in the UBT repository (2026).

\end{thebibliography}

\end{document}
